\chapter{Especificaciones Técnicas}

\section{Tipo y categoría de cables}

Para garantizar el rendimiento requerido en los \totalPisos{} niveles del pabellón de la Escuela Profesional de Ingeniería Informática y de Sistemas de la UNSAAC, se especifica el siguiente esquema de cableado:

\begin{itemize}
    \item \textbf{Cableado Horizontal (Datos, Telefonía IP y CCTV):} Se utilizará cable de par trenzado de cuatro pares del tipo \textbf{F/UTP (Foiled/Unshielded Twisted Pair) Categoría 6A} LSZH (Low Smoke Zero Halogen), marca \textbf{Panduit} modelo \texttt{PUFL6X04WH}. El blindaje general foil protege contra la diafonía exógena (Alien Crosstalk) y las interferencias electromagnéticas (EMI) típicas de los entornos de laboratorios informáticos y equipos eléctricos del edificio. Soporta velocidades de hasta 10 Gbps a frecuencias de 500 MHz, cumpliendo con el estándar IEEE 802.3an (10GBASE-T).
    \item \textbf{Cableado Vertical o Backbone:} Interconexión del Distribuidor Principal de Cableado (MDF) situado en el segundo piso con los Distribuidores de Piso (IDF) del primer, tercer y cuarto piso mediante \textbf{Fibra Óptica Multimodo OM4} de 4 hilos, con chaqueta LSZH, marca \textbf{CommScope} serie \texttt{TeraSPEED}, optimizada para transmisión láser de 10/40/100 Gbps. Se instalarán tres enlaces independientes: MDF (2do piso) a IDF (1er piso), MDF (2do piso) a IDF (3er piso) y MDF (2do piso) a IDF (4to piso), con longitudes estimadas de 20, 20 y 40 metros respectivamente.
    \item \textbf{Cableado para CCTV (Cámaras IP):} Se utilizará cable \textbf{F/UTP Cat 6A} LSZH (mismo tipo que el cableado horizontal) para las cámaras IP de videovigilancia, alimentadas mediante PoE+ desde los switches de acceso. Los puntos de cámara compartirán el mismo esquema de cableado y certificación que los puntos de datos.
    \item \textbf{Patch Cords:} Cables de parcheo F/UTP Cat 6A blindados de 2.1 m y 1.0 m, marca \textbf{Panduit} modelo \texttt{STP6X2MBU} y \texttt{STP6X1MBU}, con conectores blindados RJ-45 moldeados y alivio de tensión integrado. Se suministrarán dos patch cords por cada punto de red instalado (uno para el lado de usuario y otro en el rack).
\end{itemize}

\subsection*{Definición y clasificación de tomas y puertos de red}

Para efectos de control métrico y técnico de este expediente, se establecen las siguientes definiciones de manera unificada:
\begin{itemize}
    \item \textbf{Puerto físico (RJ-45):} Interfaz terminal hembra de 8 hilos blindada que aloja un conector RJ-45.
    \item \textbf{Punto de red (o enlace horizontal):} El canal completo que va desde el jack RJ-45 en el área de trabajo del usuario, discurre por la canalización horizontal y finaliza en el patch panel modular correspondiente en el rack. Cada punto de red equivale exactamente a un puerto activo en el switch.
    \item \textbf{Punto simple:} Toma de red compuesta por un único jack RJ-45 montado en un faceplate de 1 o 2 puertos. Se utiliza para dispositivos dedicados de un solo puerto como APs, cámaras de videovigilancia, proyectores o pizarras digitales.
    \item \textbf{Punto doble:} Toma de red compuesta por dos jacks RJ-45 independientes montados sobre el mismo faceplate de 2 puertos, conectados mediante dos cables F/UTP Cat 6A independientes al patch panel. Se emplea en puestos de trabajo (estaciones de docentes, administrativos y laboratorios) para permitir la conexión paralela de voz (teléfono IP) y datos (PC) sin conmutación local.
    \item \textbf{Punto de reserva:} Enlace de cableado horizontal completo instalado físicamente pero no conectorizado al switch de forma activa, reservado para expansiones o redundancia operativa.
\end{itemize}


\section{Normas de instalación, etiquetado y terminación}

Toda la infraestructura instalada debe ceñirse estrictamente a los estándares internacionales de telecomunicaciones. El aseguramiento de la calidad en la transmisión de datos a 10 Gbps depende no solo de la calidad de los materiales, sino fundamentalmente de la estricta adherencia a los procedimientos técnicos de instalación:

\begin{itemize}
    \item \textbf{Estándares de Diseño e Instalación:} ANSI/TIA-568.0-D (Cableado Genérico para Telecomunicaciones en las Instalaciones del Cliente) y ANSI/TIA-568.1-D (Estándar de Infraestructura de Telecomunicaciones para Edificios Comerciales). Adicionalmente, se cumple con el \textbf{Reglamento Nacional de Edificaciones (RNE)} del Perú, norma EM.010 (Redes e Infraestructura de Telecomunicaciones) y el \textbf{Código Nacional de Electricidad - Utilización (CNE)} para canalizaciones y seguridad eléctrica.
    \item \textbf{Esquema de Terminación:} Se implementará el patrón de conexión \textbf{T568B} de manera uniforme en todos los módulos RJ-45 (jacks), patch panels y patch cords de la obra.
    \item \textbf{Radio de Curvatura Mínimo:} Durante el tendido, se respetará un radio de curvatura mínimo de \textbf{4 veces el diámetro exterior del cable} (aproximadamente 32 mm para cable Cat 6A de 8 mm de diámetro). En ningún caso se excederá este límite, ya que deformaciones en la geometría de los pares trenzados provocan pérdida de retorno y degradación de la señal a 500 MHz.
    \item \textbf{Tensión Máxima de Tracción:} La fuerza de tiro durante el tendido no superará las \textbf{25 libras-fuerza (lbf)} equivalentes a 111 Newtons. Se utilizarán lubricantes apropiados para cableado estructurado en tramos de alta fricción y se evitarán tirones bruscos que puedan estirar el cable y alterar sus propiedades eléctricas.
    \item \textbf{Destrenzado Máximo:} Al momento de la terminación en los conectores (jacks y patch panels), el destrenzado de los pares no excederá los \textbf{13 mm (0.5 pulgadas)} para mantener la inmunidad al ruido (NEXT y FEXT). Valores superiores degradan irreversiblemente el desempeño del enlace.
    \item \textbf{Administración y Etiquetado:} Bajo la norma \textbf{ANSI/TIA-606-C}. Cada punto de red, cable, patch panel y gabinete poseerá una etiqueta impresa térmicamente con un código alfanumérico estructurado siguiendo el formato \textit{Piso - Cuarto de Telecomunicaciones - Rack - Panel - Puerto} (\textit{Ejemplo: 02P-MDF-A1-01} para el Segundo Piso, MDF, Rack A, Panel 1, Puerto 01). Las tomas de usuario contarán con iconos de color para diferenciar redes de datos (azul), telefonía IP (verde) y cámaras de seguridad (rojo).
\end{itemize}

\subsection*{Control de Alien Crosstalk (AXT)}

Dado que el cableado F/UTP Cat 6A opera a frecuencias de hasta 500 MHz, es crítico controlar la diafonía exógena (Alien Crosstalk) entre cables adyacentes, especialmente en los mazos densos que salen de los gabinetes MDF e IDF hacia los laboratorios y aulas:

\begin{itemize}
    \item \textbf{Sujeción:} Se prohíbe terminantemente el uso de cintillos plásticos (zip ties) apretados que deformen la geometría del cable. Se utilizará exclusivamente \textbf{velcro} (cinta enrollable) de forma holgada para agrupar los cables, permitiendo el movimiento natural y evitando puntos de presión que alteren la impedancia del cable.
    \item \textbf{Factor de Llenado:} Se limitará la ocupación de las canaletas al \textbf{40\%} de la sección transversal, conforme a la norma ANSI/TIA-569-D. Esto permite que los cables se acomoden aleatoriamente (random lay), reduciendo naturalmente la alineación de pares y cancelando el ruido entre cables vecinos.
    \item \textbf{Separación de Servicios:} Se mantendrá una separación física mínima de \textbf{200 mm (8 pulgadas)} entre el cableado de datos F/UTP Cat 6A y cualquier línea eléctrica de potencia (220V) que corra en paralelo. En áreas de espacio reducido donde esta separación no sea viable, se instalarán barreras físicas metálicas conectadas a tierra.
    \item \textbf{Cruces Perpendiculares:} Cualquier intersección necesaria entre la ruta de datos y los circuitos eléctricos se ejecutará obligatoriamente en un ángulo de \textbf{90 grados} para anular la inducción magnética.
\end{itemize}

\section{Estándares de racks, gabinetes y salas de telecomunicaciones}

La infraestructura de soporte físico se regirá bajo el estándar EIA/ECA-310-E (dimensiones de 19 pulgadas) y ANSI/TIA-569-D (espacios de telecomunicaciones). Se priorizará el uso de gabinetes cerrados (enclosures) sobre racks abiertos para garantizar la seguridad física, la protección contra polvo y humedad, y el control de acceso:

\begin{itemize}
    \item \textbf{Distribuidor de Piso (IDF - 1er Piso):} Gabinete de piso/pared de \textbf{24U}, ubicado en el cuarto de telecomunicaciones acondicionado bajo la escalera principal del primer piso. Permite alimentar las aulas y oficinas administrativas de dicho nivel sin exceder los límites de longitud horizontal del cableado.
    \item \textbf{Distribuidor Principal (MDF - 2do Piso):} Gabinete de piso autocualificado de \textbf{19 pulgadas} de ancho útil y una altura de \textbf{42U}. Profundidad mínima de 1000 mm para alojar equipos de gran profundidad y gestionar holgadamente el cableado Cat 6A. Cumplirá con el estándar EIA-310-E, con puertas microperforadas (frontal y trasera) para optimizar el flujo de aire, cerradura de seguridad con llave y paneles laterales desmontables. Alojará los servidores locales, switches troncales, patch panels del segundo piso y el panel de fibra óptica principal.
    \item \textbf{Distribuidor de Piso (IDF - 3er Piso):} Gabinete de piso/pared de \textbf{24U}, ubicado estratégicamente en un cuarto de telecomunicaciones dedicado en el tercer piso para no superar el límite de 90 metros de canalización horizontal hacia los laboratorios de cómputo.
    \item \textbf{Distribuidor de Piso (IDF - 4to Piso):} Gabinete de piso/pared de \textbf{24U}, ubicado para alimentar los cubículos de docentes, auditorio, laboratorio de investigación y control de cámaras del cuarto piso.
    \item \textbf{Ordenadores de Cable:} Se instalarán organizadores horizontales de 1U con ducto ranurado frontal y tapa desmontable entre cada patch panel, y organizadores verticales dentro de los gabinetes para asegurar un radio de curvatura correcto en los cables de parcheo y mantener el orden.
    \item \textbf{Anclaje Antisísmico:} Considerando la ubicación del proyecto en Cusco (zona sísmica), todos los gabinetes se fijarán al piso mediante pernos de expansión tipo Hilti de 3/8" o 1/2". En el caso de gabinetes adosados a pared, se verificará que el muro sea de concreto o ladrillo macizo para soportar el torque de carga. En el cuarto piso, al ser muros de drywall, se aplicará un refuerzo estructural con planchas de madera contrachapada y anclaje superior de seguridad a la losa de techo.
\end{itemize}

\subsection*{Criterios de diseño de salas de telecomunicaciones}

Cada cuarto de telecomunicaciones (MDF e IDF) deberá cumplir con los siguientes requisitos mínimos:

\begin{itemize}
    \item \textbf{Áreas de Trabajo (Clearance):} Espacio libre mínimo de \textbf{1.0 metro} tanto en la parte frontal como posterior del gabinete para permitir la apertura total de puertas, el ingreso de aire frío para ventilación y la maniobra técnica segura durante la certificación y el parcheo.
    \item \textbf{Ventilación y Climatización:} Los recintos designados requieren ventilación asistida (extractores o aire acondicionado) para disipar el calor generado por los equipos activos (switches PoE+ que pueden disipar hasta 300W cada uno) y las baterías de los UPS, evitando la degradación de la vida útil de los componentes electrónicos.
    \item \textbf{Iluminación:} Nivel de iluminación mínimo de 500 lux en el plano vertical frontal de los gabinetes para facilitar las labores de parcheo, mantenimiento y certificación.
    \item \textbf{Capacidad Portante del Piso:} Se estima que un gabinete completamente equipado (switches + UPS + baterías + cableado) puede alcanzar un peso dinámico de 150 a 250 kg. Las losas de concreto deben ser aptas para soportar esta carga puntual.
    \item \textbf{Tomas Eléctricas Dedicadas:} Cada cuarto de telecomunicaciones contará con circuitos eléctricos independientes de 20A o 30A con tierra aislada, dedicados exclusivamente a los equipos de red, minimizando el riesgo de sobrecargas por conexión de equipos no autorizados.
\end{itemize}

\subsection*{Equipos Pasivos}

\begin{itemize}
    \item \textbf{Patch Panels:} Módulos modulares descargados de 24 puertos Cat 6A angulares, marca \textbf{CommScope} modelo \texttt{760237040}, para optimizar el espacio y evitar el estrés mecánico en los patch cords. Se instalarán un total de 15 paneles distribuidos para atender de forma ordenada la densidad de puntos: 2 paneles en el IDF del 1er piso, 3 en el MDF (2do piso), 6 en el IDF del 3er piso y 4 en el IDF del 4to piso (dando un total de 360 puertos disponibles para los \totalPuntos{} puntos activos).
    \item \textbf{Jacks y Conectores:} Módulos RJ-45 Mini-Com Cat 6A blindados, marca \textbf{Panduit} modelo \texttt{CJS6X88TGY}, con tecnología de terminación por desplazamiento de aislante (IDC) y contacto de 50 micras de oro.
    \item \textbf{Faceplates y Cajas de Montaje:} Placas faceplate de 2 puertos color blanco, marca \textbf{Panduit} modelo \texttt{CFPL2EIWH}, con etiquetado iconográfico integrado para datos (azul), voz (verde) y CCTV (rojo). Se instalarán sobre cajas superficiales de PVC de 1 módulo en paredes y canaletas.
    \item \textbf{Patch Cords:} Cables de parcheo F/UTP Cat 6A blindados, marca \textbf{Panduit} modelo \texttt{STP6X2MBU} (2.1 m) y \texttt{STP6X1MBU} (1.0 m), con conectores blindados RJ-45 moldeados y alivio de tensión. Se suministrarán dos unidades por cada punto de red (uno para el lado del usuario y otro para el lado del patch panel) totalizando \totalPuntos{} unidades de cada longitud.
    \item \textbf{Organizadores de Cable:} Organizadores horizontales de 1U con ducto ranurado frontal y tapa desmontable, marca \textbf{Tripp Lite} modelo \texttt{NR2600}, instalados entre cada patch panel. Organizadores verticales dentro de cada gabinete para el guiado de cables de parcheo.
    \item \textbf{Canalizaciones:} Escalerillas portacables de tipo malla electrozincada en pasadizos principales (\metrosEscalerilla{} metros estimados), canaletas de PVC de doble vía con división física para datos y energía en oficinas y laboratorios (\metrosCanaleta{} metros), y tubería corrugada de 1 pulgada para la protección del backbone de fibra óptica entre pisos (\metrosTuboConduit{} metros).
    \item \textbf{Sistema de Puesta a Tierra:} Barra de Puesta a Tierra Principal de Telecomunicaciones (TMGB) en el MDF del segundo piso, marca \textbf{Panduit} modelo \texttt{TMGB-060}, conectada directamente al pozo a tierra general del edificio mediante conductor de cobre desnudo calibre \textbf{2 AWG} (35 mm$^2$) en el menor recorrido posible y sin empalmes. Barras secundarias (TGB) marca \textbf{Panduit} modelo \texttt{TGB-040} en los IDF del primer, tercer y cuarto piso, interconectadas a la TMGB mediante conductores de cobre aislado verde 6 AWG.
\end{itemize}

\subsection*{Equipos Activos}

\begin{itemize}
    \item \textbf{Switches de Acceso:} Switches gestionables de Capa 2/3 con soporte \textbf{PoE+ (Power over Ethernet - IEEE 802.3at)} marca \textbf{Cisco} serie \textbf{Catalyst 9200}, modelos \texttt{C9200-48P-E} (\totalSwitchesCuarentaYOcho{} unidades) y \texttt{C9200-24P-E} (\totalSwitchesVeinticuatro{} unidad), con puertos de bajada a 1 Gbps Base-T y uplinks de 10 Gbps SFP+ para el backbone de fibra óptica. Se distribuirán: 1x C9200-48P-E en el IDF del 1er piso, 1x C9200-48P-E y 1x C9200-24P-E en el MDF (2do piso), 3x C9200-48P-E en el IDF del 3er piso y 2x C9200-48P-E en el IDF del 4to piso.
    \item \textbf{Transceptores SFP+:} Módulos SFP+ multimodo 10 GBase-SR, marca \textbf{Cisco} modelo \texttt{SFP-10G-SR}, para los enlaces de fibra entre el MDF del segundo piso y los IDF del primer, tercer y cuarto piso (6 unidades en total: 2 transceptores por enlace).
    \item \textbf{Puntos de Acceso Inalámbrico (WiFi 6):} Access Points de interior doble banda (2.4 GHz / 5 GHz) con soporte \textbf{WiFi 6 (802.11ax)}, marca \textbf{Cisco} modelo \texttt{C9120AXI}, alimentados mediante PoE+. Se instalarán estratégicamente en pasadizos y ambientes comunes de los cuatro niveles (\totalAP{} unidades en total, distribuidos a razón de 4 APs por piso) para brindar cobertura inalámbrica en laboratorios, auditorio, sala de reuniones, círculos de estudio, biblioteca y áreas administrativas.
    \item \textbf{Cámaras IP (CCTV):} Cámaras IP fijas de interior con resolución mínima de 2 MP (1920x1080), PoE+, marca \textbf{Dahua} o \textbf{Hikvision}, conectadas a los switches de acceso PoE+. Se proyecta un total de \totalCamaras{} cámaras (3 unidades por nivel) integradas con el NVR (Network Video Recorder) de 16 canales que se ubicará en el gabinete MDF del segundo piso para la grabación y gestión centralizada de video.
    \item \textbf{UPS:} Sistemas de alimentación ininterrumpida en línea de doble conversión, marca \textbf{Tripp Lite} modelo \texttt{SU3000RTX} (3kVA), un equipo por cada gabinete (IDF 1er piso + MDF 2do piso + IDF 3er piso + IDF 4to piso = \totalUPS{} unidades en total), con autonomía mínima de 15 minutos a plena carga para salvaguardar la operatividad de los switches, APs y cámaras.
    \item \textbf{PDU:} Unidades de Distribución de Energía rackeables, marca \textbf{Tripp Lite} modelo \texttt{PDU1230}, con entrada monofásica y múltiples salidas NEMA 5-15R/5-20R, una por cada gabinete (\totalPDU{} unidades en total).
\end{itemize}

\section{Canalizaciones, registros y detalles constructivos}

El sistema de canalizaciones constituye la infraestructura de soporte físico que alberga y protege el cableado estructurado. Se divide en canalización principal, canalización secundaria y sistema de registros:

\subsection*{Canalización Principal (Escalerillas Tipo Malla)}

\begin{itemize}
    \item \textbf{Tipo:} Escalerilla portacables de malla electrozincada (tipo rejilla), fabricada en acero al carbono con recubrimiento electrozincado, cumpliendo con la norma ANSI/TIA-569-D y NEMA VE 1.
    \item \textbf{Dimensiones:} Ancho mínimo de 300 mm y profundidad de 50 mm, con separación entre alambres de 50 mm x 100 mm.
    \item \textbf{Soportería:} Instalada cada 1.5 metros sobre soportes tipo L o mensualas metálicas fijadas a losas o paredes mediante pernos de expansión. En tramos horizontales en pasillos principales, se fijará al techo o bajo el falso techo.
    \item \textbf{Rutas:} Discurrirá por los pasillos principales de distribución de cada nivel (primer, segundo, tercero y cuarto piso), conectando los gabinetes MDF/IDF con las zonas de distribución secundaria.
    \item \textbf{Tapas:} En tramos donde la escalerilla cruce zonas accesibles al público, se instalará tapa metálica desmontable para protección mecánica del cableado.
\end{itemize}

\subsection*{Canalización Secundaria (Canaletas PVC y Tuberías)}

\begin{itemize}
    \item \textbf{Canaletas PVC:} Canaletas de PVC rígido autoextinguible (libre de halógenos), de doble compartimiento con divisor físico central para segregación datos/energía. Dimensiones mínimas de 40 mm x 20 mm (2 vías). Se instalarán adosadas a pared en oficinas, laboratorios y cubículos, desde la escalerilla principal hasta cada punto de red.
    \item \textbf{Tubería Conduit:} Para el backbone de fibra óptica entre pisos se utilizará tubería conduit metálica (EMT) de 1 pulgada de diámetro, con accesorios de unión y curvas de radio largo que garanticen el radio de curvatura mínimo de la fibra OM4 (10 veces el diámetro exterior).
    \item \textbf{Factor de Llenado:} En ningún caso la ocupación de cables superará el \textbf{40\%} del área transversal de la canaleta o tubería, conforme a la norma ANSI/TIA-569-D.
    \item \textbf{Separación de Servicios:} Se mantendrá una separación mínima de 200 mm entre canalizaciones de datos y circuitos eléctricos de 220V. En espacios reducidos se usarán canaletas con divisor metálico conectado a tierra.
\end{itemize}

\subsection*{Registros de Paso y Cajas de Derivación}

\begin{itemize}
    \item \textbf{Registros de Paso:} Se instalarán cajas de registro metálicas o de PVC cada 15 metros lineales de canalización, en cada cambio de dirección superior a 30 grados y en los puntos de derivación hacia ambientes específicos. Dimensiones mínimas de 100 mm x 100 mm x 50 mm.
    \item \textbf{Cajas de Paso en Pisos:} Para los tendidos verticales entre niveles (montantes), se instalarán buzones de paso metálicos en cada losa, con dimensiones mínimas de 200 mm x 200 mm, provistos de sellos cortafuego intumescentes para mantener la integridad de la losa contra incendios.
    \item \textbf{Sellado Cortafuego:} Todo pase de cableado entre pisos o a través de muros cortafuego deberá ser sellado con masilla intumescente certificada, manteniendo la clasificación de resistencia al fuego de la estructura (mínimo 2 horas RF-120).
\end{itemize}

\subsection*{Detalles Constructivos}

\begin{itemize}
    \item \textbf{Soportería de Bandejas:} Las escalerillas se soportarán sobre estructuras metálicas tipo L cada 1.5 m, fijadas a losa o pared con pernos de expansión de 3/8". En tramos paralelos a tuberías eléctricas, se interpondrá una barrera metálica conectada a tierra.
    \item \textbf{Curvas y Derivaciones:} Se utilizarán curvas prefabricadas de radio largo (mínimo 300 mm) para garantizar el radio de curvatura del cable Cat 6A. No se aceptarán curvas realizadas en obra mediante doblado de canaletas.
    \item \textbf{Fijación de Canaletas:} Las canaletas PVC se fijarán a la pared cada 40 cm mediante tacos y tornillos de expansión, con tapas de registro cada 2 metros para facilitar el mantenimiento.
    \item \textbf{Amarre de Cables:} Los cables se fijarán a las escalerillas mediante velcro cada 1 metro, sin apretar para no deformar la geometría del par trenzado. Se prohíbe el uso de cintillos plásticos.
\end{itemize}

\section{Requisitos eléctricos y de puesta a tierra}

\subsection*{Sistema de Puesta a Tierra para Telecomunicaciones}

En cumplimiento estricto con las normas \textbf{ANSI/TIA-607-D} y \textbf{ANSI/TIA-607-C}, se implementará un sistema de puesta a tierra dedicado para telecomunicaciones que garantice la protección de los equipos activos y la estabilidad de la transmisión de datos:

\begin{itemize}
    \item \textbf{Barra Principal (TMGB):} Ubicada en el MDF del segundo piso, marca \textbf{Panduit} modelo \texttt{TMGB-060}, conectada directamente al pozo a tierra general del edificio mediante conductor de cobre desnudo calibre \textbf{2 AWG} (35 mm$^2$) en el menor recorrido posible y sin empalmes. Esta barra servirá como punto de referencia de tierra para todo el sistema de telecomunicaciones de los \totalPisos{} niveles.
    \item \textbf{Barras Secundarias (TGB):} Barras secundarias (TGB) marca \textbf{Panduit} modelo \texttt{TGB-040} en los IDF del primer, tercer y cuarto piso, interconectadas a la TMGB mediante conductores de cobre aislado color verde de calibre \textbf{6 AWG} (13.3 mm$^2$) por el interior de la canalización del backbone.
    \item \textbf{Vinculación (Bonding):} Todas las estructuras metálicas pasivas se interconectarán a las barras TGB/TMGB:
        \begin{itemize}
            \item Chasis de gabinetes y racks.
            \item Bandejas portacables metálicas (escalerillas tipo malla).
            \item Blindaje de patch panels y cables F/UTP (mediante la barra de tierra del patch panel).
            \item Canaletas metálicas y tuberías EMT del backbone.
            \item PDU y UPS (mediante conductor de tierra del equipo).
        \end{itemize}
    \item \textbf{Conductor de Enlace Equipotencial:} Se tenderá un conductor de cobre aislado verde de calibre 6 AWG a lo largo de las escalerillas portacables, conectado a cada tramo de bandeja y a cada gabinete, garantizando la continuidad eléctrica de todo el sistema de soporte.
    \item \textbf{Resistencia de Puesta a Tierra:} La resistencia del sistema de puesta a tierra, medida con un telurómetro, no deberá superar los \textbf{5 ohmios} en la TMGB, conforme a la norma ANSI/TIA-607-D y el CNE (Código Nacional de Electricidad - Utilización).
\end{itemize}

\subsection*{Alimentación Eléctrica Estabilizada}

\begin{itemize}
    \item \textbf{Circuitos Dedicados:} Cada cuarto de telecomunicaciones (MDF e IDF) contará con circuitos eléctricos exclusivos de 20A o 30A, con tomas NEMA 5-20R o L5-30R con tierra aislada, independientes de los circuitos de iluminación y tomacorrientes generales.
    \item \textbf{Sistema UPS:} Los equipos activos de cada gabinete se alimentarán mediante UPS en línea de doble conversión marca \textbf{Tripp Lite} modelo \texttt{SU3000RTX} (3kVA), con autonomía mínima de 15 minutos a plena carga. Se instalarán \totalUPS{} unidades en total (uno por cada gabinete de comunicaciones).
    \item \textbf{Generador de Respaldo:} Se propone e implementa un generador diésel de emergencia de \textbf{50 kVA} con interruptor de transferencia automática (ATS) marca Cummins, para alimentar de forma ininterrumpida los circuitos dedicados de los cuatro cuartos de telecomunicaciones y el sistema de climatización del MDF en caso de fallas de la red eléctrica comercial.
\end{itemize}

\subsection*{Diagrama de Puesta a Tierra}

Se incluirá en los planos un diagrama unifilar del sistema de puesta a tierra de telecomunicaciones que muestre:

\begin{itemize}
    \item Ubicación de la TMGB en el MDF (2do piso) con su conexión al pozo a tierra del edificio.
    \item Ubicación de las TGB en los IDF (1er, 3er y 4to piso).
    \item Trayectoria y calibre de los conductores de enlace (bonding) entre barras.
    \item Identificación de todos los elementos metálicos a conectar (racks, bandejas, canaletas, patch panels).
    \item Detalle de conexión del blindaje de los cables F/UTP a la barra de tierra del patch panel.
    \item Especificación de la resistencia máxima admisible (5 ohmios) y método de medición.
\end{itemize}

\section{Procedimientos de certificación}

Una vez finalizado el tendido y terminación del total de \textbf{\totalPuntos{} puntos de red}, se procederá a la validación técnica total del sistema:

\subsection*{Alcance y Metodología}

\begin{itemize}
    \item \textbf{Cobertura:} Se someterá a prueba el 100\% de los puntos de red de cobre instalados en los \totalPisos{} niveles del pabellón.
    \item \textbf{Configuración de Prueba:} Las pruebas se realizarán bajo la configuración de \textbf{Enlace Permanente (Permanent Link)}, evaluando el tramo desde el Patch Panel en el gabinete hasta la toma (Jack RJ-45) en el área de trabajo. No se utilizará la configuración de Canal (Channel) para la certificación de instalación, ya que esta depende de los patch cords del usuario final.
    \item \textbf{Instrumentación:} Se utilizará un certificador de cableado de campo de Nivel IV o superior (\textbf{Fluke Networks DSX-8000} calibrado), con certificado de calibración vigente emitido por el fabricante o laboratorio autorizado con una antigüedad no mayor a 12 meses.
\end{itemize}

\subsection*{Parámetros de Evaluación}

El procedimiento de Autotest verificará los siguientes parámetros críticos para la operación a 500 MHz (10GBASE-T) según la norma ANSI/TIA-568.2-D:

\begin{itemize}
    \item \textbf{Mapa de Cableado (Wire Map):} Verificación de continuidad pin a pin, cortocircuitos, pares abiertos, pares cruzados y pares divididos (split pairs). Confirmación de la correcta terminación T568B y continuidad del blindaje (shield) general del cable F/UTP.
    \item \textbf{Longitud (Length):} Validación de que el enlace permanente no exceda los 90 metros. Cualquier enlace superior será considerado FALLA automática.
    \item \textbf{Pérdida de Inserción (Insertion Loss):} Medición de la atenuación de la señal a lo largo del cable.
    \item \textbf{Pérdida de Retorno (Return Loss):} Medición de la energía reflejada por diferencias de impedancia (curvas excesivas o malas terminaciones). Parámetro crítico para 10 Gbps.
    \item \textbf{Diafonía NEXT (Near-End Crosstalk):} Interferencia de un par sobre otro en el extremo cercano.
    \item \textbf{PS NEXT (Power Sum NEXT):} Efecto combinado de tres pares interfiriendo sobre el cuarto par. Crítico en Cat 6A para operación 10GBASE-T sobre 4 pares simultáneos.
    \item \textbf{ACR-N y ACR-F:} Relación Atenuación-Diafonía en extremos cercano y lejano. Indicadores de la calidad del canal de transmisión.
    \item \textbf{Delay Skew:} Diferencia de tiempo de propagación entre pares. No debe exceder los 50 ns para garantizar la integridad de la transmisión 10GBASE-T.
\end{itemize}

*(Para detalles referentes a la ejecución en obra, muestreo por el supervisor, protocolos de corrección de fallas y formatos del dossier de calidad, consúltese el Capítulo 6 del presente expediente).*

\subsection*{Pruebas de Fibra Óptica (Backbone)}

La certificación de la fibra óptica multimodo OM4 se realizará bajo los parámetros del estándar ANSI/TIA-568.3-D:
\begin{itemize}
    \item \textbf{Certificación de Nivel 1 (Tier 1):} Medición de la pérdida de inserción (atenuación) y la longitud mediante un equipo LSPM (Fuente de Luz y Medidor de Potencia), evaluando ambas direcciones y en las longitudes de onda de 850 nm y 1300 nm para cada uno de los hilos de fibra OM4 de los tres enlaces de backbone.
    \item \textbf{Límites de Desempeño:} La atenuación total de cada enlace no superará el presupuesto de pérdida calculado (típicamente $\le 2.0\text{ dB}$ para enlaces de fibra óptica cortos de hasta 40 metros, incluyendo acopladores y empalmes).
\end{itemize}

\section{Criterios de compatibilidad electromagnética}

Para garantizar la inmunidad del sistema de cableado estructurado a las interferencias electromagnéticas presentes en el entorno del pabellón, se aplicarán los siguientes criterios:

\begin{itemize}
    \item \textbf{Clasificación del Entorno:} El entorno electromagnético del pabellón se clasifica como \textbf{MICE 1} (Entorno de Oficina/Comercial) según la norma ISO/IEC 11801, lo que permite el uso de cableado F/UTP sin requerir apantallamiento adicional siempre que se respeten las distancias de separación.
    \item \textbf{Fuentes de Interferencia Identificadas:} Iluminación fluorescente/LED en falsos techos, tableros eléctricos de distribución en zonas comunes (tercer piso), motores eléctricos de sistemas de ventilación y cableado eléctrico de 220V en laboratorios.
    \item \textbf{Mitigación:} Aplicación estricta de las distancias de separación ANSI/TIA-569-D (200 mm de paralelismo con circuitos eléctricos, cruces a 90 grados), uso de cable blindado F/UTP y conexión del blindaje a tierra en ambos extremos mediante patch panels con barra de tierra integrada.
\end{itemize}
